% GENERATED FILE. Edit canonical lessons, then run npm run generate:books.
% canonical-source-hash: fnv1a64:21765b0a9aff5dbd
% canonical-lessons: SA-C382-hitam, SA-C382-kalyanam, SA-C382-sreyah, SA-C382-visramyati, SA-C382-satkaroti

\chapter{For Our Good}
\label{ch:sa-for-our-good}
\hlchaptermodality{382}

\begin{chapteropening}
\textbf{By the end of this chapter:} \emph{I can say what would be good for someone.}

Complete the last lesson of chapter 382: I can say what would be good for someone.
\end{chapteropening}

\section[\texorpdfstring{\sk{हितम्} (hitam)}{हितम् (hitam)}]{\sk{हितम्} (hitam) — benefit; welfare}

\label{lesson:SA-C382-hitam}



\begin{quote}
Before the new one: say the Sanskrit for repeatedly, then the Sanskrit for sums up.
\end{quote}

\subsection*{You'll want to know: \sk{हितम्}}
\textbf{\sk{हितम्}} — \emph{hitam} — \textquotedblleft{}benefit; welfare\textquotedblright{}.

\begin{quote}
\sk{यदि} \sk{एतत्} \sk{मम} \sk{हितम्}, \sk{तर्हि} \sk{अहम्} \sk{अवश्यं} \sk{करिष्यामि।}
\end{quote}

\begin{quote}
\emph{yadi etat mama hitam, tarhi aham avaśyaṃ kariṣyāmi.}
\end{quote}

\begin{quote}
If this is for my good, I will do it by all means.
\end{quote}

A neuter noun.

\begin{cousinweb}
\textbf{\sk{हितम्}} (\emph{hitam}) comes from the root √dhā, \textquotedblleft{}to put\textquotedblright{}, whose participle MW gives as \emph{hita}, \textquotedblleft{}put, placed\textquotedblright{}; MW gives \emph{hita} the senses \textquotedblleft{}beneficial, advantageous, salutary\textquotedblright{} (MW).
\end{cousinweb}

\begin{grammarlens}[title={What you've built}]
One more word for saying what happens if.
\end{grammarlens}

\subsection*{Your turn}
\begin{itemize}
\raggedright
  \item \emph{Say it:} \emph{hitam}
  \item \emph{Say it:} \emph{hitam}, once more
  \item \emph{Say it:} \emph{yadi etat mama hitam, tarhi aham avaśyaṃ kariṣyāmi.} — If this is for my good, I will do it by all means.
  \item \emph{Recall:} say the Sanskrit for repeatedly, then the Sanskrit for sums up, then say \emph{hitam} again
\end{itemize}

\begin{tcolorbox}[breakable,colback=teal!4,colframe=teal!35!black,title={Before you move on}]
What does \sk{हितम्} mean? (\textquotedblleft{}Benefit; welfare\textquotedblright{}.) Say it once more.
\end{tcolorbox}
\section[\texorpdfstring{\sk{कल्याणम्} (kalyāṇam)}{कल्याणम् (kalyāṇam)}]{\sk{कल्याणम्} (kalyāṇam) — welfare; good fortune}

\label{lesson:SA-C382-kalyanam}



\begin{quote}
Before the new one: say the Sanskrit for sums up, then the Sanskrit for benefit.
\end{quote}

\subsection*{You'll want to know: \sk{कल्याणम्}}
\textbf{\sk{कल्याणम्}} — \emph{kalyāṇam} — \textquotedblleft{}welfare; good fortune\textquotedblright{}.

\begin{quote}
\sk{यदि} \sk{सः} \sk{सत्यं} \sk{वदति}, \sk{तर्हि} \sk{तस्य} \sk{कल्याणं} \sk{भविष्यति।}
\end{quote}

\begin{quote}
\emph{yadi saḥ satyaṃ vadati, tarhi tasya kalyāṇaṃ bhaviṣyati.}
\end{quote}

\begin{quote}
If he tells the truth, good will come to him.
\end{quote}

A neuter noun.

\begin{cousinweb}
\textbf{\sk{कल्याणम्}} (\emph{kalyāṇam}) MW gives \emph{kalyāṇa} \textquotedblleft{}happy, prosperous, fortunate\textquotedblright{} and the neuter \textquotedblleft{}good fortune, happiness, prosperity\textquotedblright{}; MW gives no derivation, so its origin is uncertain (MW).
\end{cousinweb}

\begin{grammarlens}[title={What you've built}]
One more word for saying what happens if.
\end{grammarlens}

\subsection*{Your turn}
\begin{itemize}
\raggedright
  \item \emph{Say it:} \emph{kalyāṇam}
  \item \emph{Say it:} \emph{kalyāṇam}, once more
  \item \emph{Say it:} \emph{yadi saḥ satyaṃ vadati, tarhi tasya kalyāṇaṃ bhaviṣyati.} — If he tells the truth, good will come to him.
  \item \emph{Recall:} say the Sanskrit for sums up, then the Sanskrit for benefit, then say \emph{kalyāṇam} again
\end{itemize}

\begin{tcolorbox}[breakable,colback=teal!4,colframe=teal!35!black,title={Before you move on}]
What does \sk{कल्याणम्} mean? (\textquotedblleft{}Welfare; good fortune\textquotedblright{}.) Say it once more.
\end{tcolorbox}
\section[\texorpdfstring{\sk{श्रेयः} (śreyaḥ)}{श्रेयः (śreyaḥ)}]{\sk{श्रेयः} (śreyaḥ) — the better course; welfare}

\label{lesson:SA-C382-sreyah}



\begin{quote}
Before the new one: say the Sanskrit for benefit, then the Sanskrit for welfare.
\end{quote}

\subsection*{You'll want to know: \sk{श्रेयः}}
\textbf{\sk{श्रेयः}} — \emph{śreyaḥ} — \textquotedblleft{}the better course; welfare\textquotedblright{}.

\begin{quote}
\sk{यदि} \sk{भवान्} \sk{रुग्णः}, \sk{तर्हि} \sk{विश्रामः} \sk{श्रेयः।}
\end{quote}

\begin{quote}
\emph{yadi bhavān rugṇaḥ, tarhi viśrāmaḥ śreyaḥ.}
\end{quote}

\begin{quote}
If you are ill, rest is the better course.
\end{quote}

A neuter noun.

\begin{cousinweb}
\textbf{\sk{श्रेयः}} (\emph{śreyaḥ}) MW gives \emph{śreyas} \textquotedblleft{}better\textquotedblright{} and, as a neuter noun, \textquotedblleft{}the better state, the better fortune\textquotedblright{}; MW gives no derivation here, so its origin is uncertain (MW).
\end{cousinweb}

\begin{grammarlens}[title={What you've built}]
One more word for saying what happens if.
\end{grammarlens}

\subsection*{Your turn}
\begin{itemize}
\raggedright
  \item \emph{Say it:} \emph{śreyaḥ}
  \item \emph{Say it:} \emph{śreyaḥ}, once more
  \item \emph{Say it:} \emph{yadi bhavān rugṇaḥ, tarhi viśrāmaḥ śreyaḥ.} — If you are ill, rest is the better course.
  \item \emph{Recall:} say the Sanskrit for benefit, then the Sanskrit for welfare, then say \emph{śreyaḥ} again
\end{itemize}

\begin{tcolorbox}[breakable,colback=teal!4,colframe=teal!35!black,title={Before you move on}]
What does \sk{श्रेयः} mean? (\textquotedblleft{}The better course; welfare\textquotedblright{}.) Say it once more.
\end{tcolorbox}
\section[\texorpdfstring{\sk{विश्राम्यति} (viśrāmyati)}{विश्राम्यति (viśrāmyati)}]{\sk{विश्राम्यति} (viśrāmyati) — rests}

\label{lesson:SA-C382-visramyati}



\begin{quote}
Before the new one: say the Sanskrit for welfare, then the Sanskrit for the better course.
\end{quote}

\subsection*{You'll want to know: \sk{विश्राम्यति}}
\textbf{\sk{विश्राम्यति}} — \emph{viśrāmyati} — \textquotedblleft{}rests\textquotedblright{}.

\begin{quote}
\sk{यदि} \sk{भवान्} \sk{श्रान्तः}, \sk{तर्हि} \sk{क्षणं} \sk{विश्राम्यतु।}
\end{quote}

\begin{quote}
\emph{yadi bhavān śrāntaḥ, tarhi kṣaṇaṃ viśrāmyatu.}
\end{quote}

\begin{quote}
If you are tired, rest for a moment.
\end{quote}

Past (he …): \sk{सः} \sk{वृक्षस्य} \sk{छायायां} \sk{विश्रान्तवान्}, \emph{saḥ vṛkṣasya chāyāyāṃ viśrāntavān}.

\begin{cousinweb}
\textbf{\sk{विश्राम्यति}} (\emph{viśrāmyati}) is built on \emph{vi} and the root √śram, \textquotedblleft{}to be weary\textquotedblright{}; MW gives \emph{vi-√śram} \textquotedblleft{}to rest, to repose\textquotedblright{} and \emph{vi-śrānta} \textquotedblleft{}rested\textquotedblright{} (MW).
\end{cousinweb}

\begin{grammarlens}[title={What you've built}]
One more word for saying what happens if.
\end{grammarlens}

\subsection*{Your turn}
\begin{itemize}
\raggedright
  \item \emph{Say it:} \emph{viśrāmyati}
  \item \emph{Say it:} \emph{viśrāmyati}, once more
  \item \emph{Say it:} \emph{yadi bhavān śrāntaḥ, tarhi kṣaṇaṃ viśrāmyatu.} — If you are tired, rest for a moment.
  \item \emph{Recall:} say the Sanskrit for welfare, then the Sanskrit for the better course, then say \emph{viśrāmyati} again
\end{itemize}

\begin{tcolorbox}[breakable,colback=teal!4,colframe=teal!35!black,title={Before you move on}]
What does \sk{विश्राम्यति} mean? (\textquotedblleft{}Rests\textquotedblright{}.) Say it once more.
\end{tcolorbox}
\section[\texorpdfstring{\sk{सत्करोति} (satkaroti)}{सत्करोति (satkaroti)}]{\sk{सत्करोति} (satkaroti) — receives with honour}

\label{lesson:SA-C382-satkaroti}



\begin{quote}
Before the new one: say the Sanskrit for the better course, then the Sanskrit for rests.
\end{quote}

\subsection*{You'll want to know: \sk{सत्करोति}}
\textbf{\sk{सत्करोति}} — \emph{satkaroti} — \textquotedblleft{}receives with honour\textquotedblright{}.

\begin{quote}
\sk{यदि} \sk{अतिथिः} \sk{आगच्छति}, \sk{तर्हि} \sk{माता} \sk{तं} \sk{सत्करोति।}
\end{quote}

\begin{quote}
\emph{yadi atithiḥ āgacchati, tarhi mātā taṃ satkaroti.}
\end{quote}

\begin{quote}
If a guest comes, mother receives him with honour.
\end{quote}

Past (he …): \sk{सः} \sk{अतिथिं} \sk{सत्कृतवान्}, \emph{saḥ atithiṃ satkṛtavān}.

\begin{cousinweb}
\textbf{\sk{सत्करोति}} (\emph{satkaroti}) is built from \emph{sat}, \textquotedblleft{}good\textquotedblright{}, and the root √kṛ, \textquotedblleft{}to do\textquotedblright{}; MW gives \emph{sat-√kṛ} \textquotedblleft{}to treat well or with respect, to honour, to receive hospitably\textquotedblright{} (MW).
\end{cousinweb}

\begin{grammarlens}[title={What you've built}]
One more word for saying what happens if.
\end{grammarlens}

\subsection*{Your turn}
\begin{itemize}
\raggedright
  \item \emph{Say it:} \emph{satkaroti}
  \item \emph{Say it:} \emph{satkaroti}, once more
  \item \emph{Say it:} \emph{yadi atithiḥ āgacchati, tarhi mātā taṃ satkaroti.} — If a guest comes, mother receives him with honour.
  \item \emph{Recall:} say the Sanskrit for the better course, then the Sanskrit for rests, then say \emph{satkaroti} again
\end{itemize}

\begin{tcolorbox}[breakable,colback=teal!4,colframe=teal!35!black,title={Before you move on}]
What does \sk{सत्करोति} mean? (\textquotedblleft{}Receives with honour\textquotedblright{}.) Say it once more.
\end{tcolorbox}
